\documentclass[11pt,a4paper]{article}
\usepackage[utf8]{inputenc}
\usepackage[T1]{fontenc}
\usepackage[top=2.5cm,bottom=2.5cm,left=2.5cm,right=2.5cm]{geometry}
\usepackage{amsmath,amssymb,amsthm,mathtools}
\usepackage[numbers,sort&compress]{natbib}
\usepackage{float}
\usepackage{xr-hyper}
\usepackage{hyperref}
\makeatletter\g@addto@macro{\UrlBreaks}{\UrlOrds}\makeatother
\usepackage{enumitem}
\usepackage{booktabs}
\usepackage{tikz}
\usetikzlibrary{arrows.meta,positioning,shapes.geometric,fit,decorations.pathreplacing,calc}
\usepackage{xcolor}
\setlength{\emergencystretch}{1em}
\DeclareMathOperator*{\argmin}{arg\,min}
\DeclareMathOperator*{\argmax}{arg\,max}


\theoremstyle{plain}
\newtheorem{theorem}{Theorem}[section]
\newtheorem{proposition}[theorem]{Proposition}
\newtheorem{lemma}[theorem]{Lemma}
\theoremstyle{definition}
\newtheorem{definition}[theorem]{Definition}
\newtheorem{example}[theorem]{Example}
\theoremstyle{remark}
\newtheorem{remark}[theorem]{Remark}

\definecolor{br1}{RGB}{31,119,180}
\definecolor{br2}{RGB}{255,127,14}
\definecolor{br3}{RGB}{44,160,44}
\definecolor{br4}{RGB}{214,39,40}
\definecolor{br5}{RGB}{148,103,189}
\definecolor{br6}{RGB}{140,86,75}
\hypersetup{colorlinks=true,linkcolor=br1,citecolor=br3,urlcolor=br4}
\externaldocument[S:]{scoping}
\begin{document}

\title{\textbf{Material removed from the scoping document}}
\author{L.\ Fiaschi \\ \small ETH Z\"urich}
\date{\today}
\maketitle

This file holds the material that was removed from the scoping document \texttt{scoping.tex} on 10 November 2026.
Each section gives the reason for the removal and then the removed text, unchanged.
References to sections, equations, tables and results of the scoping document point to its labels.

\section{Compound Poisson row of the table of damage-accumulation processes}

Reason: compound Poisson is not implemented and not planned.

\begin{center}
\small
\begin{tabular}{@{}p{3cm}p{3.6cm}p{2.2cm}cp{3.5cm}@{}}
\toprule
\textbf{Model} & \textbf{Increment} $X(t)-X(s)$ & \textbf{Marginal} $X(t)$ & \textbf{Mon.} & \textbf{Closure under $+$} \\
\midrule
Compound Poisson \cite{EsaryMarshallProschan1973,Wang2022}
  & $\sum_{n=N(s)+1}^{N(t)} Y_n$
  & $\mathrm{CP}(\nu t,\,F_Y)$
  & If $Y{\ge}0$
  & Always \\
\bottomrule
\end{tabular}
\end{center}

\section{Paragraph ``Compound Poisson process''}

Reason: compound Poisson is not implemented and not planned.

\paragraph{Compound Poisson process.}
The compound Poisson process
\begin{equation}
  X(t) = \sum_{n=1}^{N(t)} Y_n
\end{equation}
models degradation by shocks~\cite{EsaryMarshallProschan1973}.
The shocks arrive according to a Poisson process $N(t)$ with rate~$\nu$.
Each shock causes a random damage~$Y_n$, and the $Y_n$ are independent draws from one distribution~$F_Y$.

When $Y_n \ge 0$, the process is monotone.
The sum of independent compound Poisson processes is always a compound Poisson process.
This model suits discrete damage events, for example impacts and overloads, rather than continuous wear.
Wang~\cite{Wang2022} gives a recent application to structural reliability.


\section{Compound Poisson row of the table of discrete-time models}

Reason: compound Poisson is not implemented and not planned.

\begin{center}
\small
\begin{tabular}{@{}p{2.1cm}p{3.4cm}p{4.0cm}p{4.2cm}@{}}
\toprule
\textbf{Model} & \textbf{Increment} $\Delta D_t$ & \textbf{Marginal} $D_k$ (homogeneous) & \textbf{Closure condition} \\
\midrule
Compound Poisson
  & $\sum_{n=1}^{N_t} Y_n$
  & $\mathrm{CP}\!\bigl(\sum_t\nu_t,\,F_Y\bigr)$
  & Always; same $F_Y$ keeps the jump law \\
\bottomrule
\end{tabular}
\end{center}

\section{Deterministic S--N curve when the scatter is negligible}

Reason: a deterministic model is not among the five models.

If the scatter is small, a deterministic S--N curve, for example the mean curve or a design percentile, may be an acceptable approximation.
The problem then reduces to the deterministic Miner's rule, and all constraints become linear.

\section{Route II rows of the table of conservative surrogates}

Reason: Route II is not part of the formulation.

\begingroup
\footnotesize
\setlength{\tabcolsep}{4.5pt}
\begin{center}
\begin{tabular}{@{}p{1.7cm}p{1.9cm}p{1.15cm}p{0.95cm}p{2.9cm}p{2.9cm}p{1.5cm}@{}}
\toprule
\textbf{Method} & \textbf{Inputs} & \textbf{Cons.?} & \textbf{$\widetilde{D}$?} & \textbf{Pros} & \textbf{Cons} & \textbf{Ref.} \\
\midrule
\multicolumn{7}{@{}l}{\textit{Route II --- dominate each increment, then convolve (explicit $\widetilde{D}$)}} \\
Moment-matched Gamma / IG
  & mean, var
  & No & Yes
  & explicit law; closed under $+$; ordering tables apply
  & matches moments but does not dominate (survival curves cross)
  & \cite{vanNoortwijk2009,YeChen2014} \\[3pt]
Dominating shared-rate Gamma
  & moments $+$ empirical survival
  & Yes & Yes
  & conservative \emph{and} explicit; serves reliability and loop; shape linear in the schedule
  & shared rate $+$ dominance margin add slack; commits to Gamma
  & \cite{ShakedShanthikumar2007} \\
\bottomrule
\end{tabular}
\end{center}
\endgroup

\section{SOS2 sentence of the Bernstein constraint}

Reason: the code does not use SOS2.

This minimum is piecewise-linear (PWL), and an equality with a special ordered set of type~2 (SOS2: a set of variables of which at most two consecutive ones are non-zero) encodes it exactly.
Both replacements imply~\eqref{S:eq:bernstein_constraint}, so they stay conservative: a fixed tangent gives one linear constraint in the counts, and the minimum of tangents gives linear constraints together with the SOS2 set.

\section{Route II surrogates}

Reason: Route II is not part of the formulation.

Route~II gives an explicit surrogate law, not a bound.
\begin{itemize}[nosep]
  \item \textbf{Moment-matched Gamma or inverse Gaussian.} Matching only the first two moments of a family closed under convolution is convenient: Gamma, with shape and rate
    \begin{equation}
      \alpha = \mathbb{E}[\Delta D]^2/\operatorname{Var}[\Delta D]
      \quad\text{and}\quad
      \beta = \mathbb{E}[\Delta D]/\operatorname{Var}[\Delta D],
    \end{equation}
    or inverse Gaussian.
    The sum then stays in the family under the closure conditions of Table~\ref{S:tab:discrete_models}.
    But this choice is \emph{not} conservative: a law with matched moments in general crosses the empirical survival function and does not dominate it.
  \item \textbf{Dominating shared-rate Gamma.} Using more of the empirical information than the moments restores conservatism.
    Fix a common rate $\beta^\star$, and, for each mission, increase the shape $\alpha_j$ until the survival function of the surrogate dominates the empirical one:
    \begin{equation}
      \overline{\mathrm{Ga}}(x;\alpha_j,\beta^\star) \ge \widehat{\bar F}_j(x) \quad \text{for all } x.
    \end{equation}
    This is possible exactly when some recorded increment is~$0$.
    If every record is positive, $\widehat{\bar F}_j(x) = 1$ below the smallest record, while $\overline{\mathrm{Ga}}(x;\alpha,\beta^\star) < 1$ for every $x > 0$, so no Gamma law dominates.
    If a share $p_0 > 0$ of the records is~$0$, then $\widehat{\bar F}_j(x) \le 1 - p_0$ for $x > 0$, and a large $\alpha_j$ dominates, because $\overline{\mathrm{Ga}}(x_{\max};\alpha,\beta^\star) \to 1$ as $\alpha \to \infty$, with $x_{\max}$ the largest record.
    The accumulated surrogate is then $\widetilde{D} = \mathrm{Ga}(A,\beta^\star)$, by the closure under a shared rate, and $D \le_{\mathrm{st}} \widetilde{D}$ by Proposition~\ref{S:prop:st_closure}.
    Its law has a cumulative shape
    \begin{equation}
      A = \sum_{t,j} x_{ijt}\,\alpha_j
    \end{equation}
    that is linear in the schedule.
    The reliability constraint reduces to the scalar condition $A \le A^\star(\varepsilon)$.
    In the absence of repairs, the loop constraint reduces to the rows $A_{H_1} \le y^{A}$ and $A_T \le y^{A}$ on the cumulative shape, with a start $y^{A}$ that is free above $A_{H_1}$.
    The cost is the slack that the shared rate and the dominance margin add.
\end{itemize}

\section{Sentence ``Route II is not part of the formulation''}

Reason: this sentence refers to the removed Route II.

Route~II is not part of the formulation.

\section{Trade-off against Route II}

Reason: the trade-off only compares the Route-I bounds with the removed Route II.

One trade-off cannot be removed.
\begin{itemize}[nosep]
  \item \textbf{Route-I bounds.} The bounds assume nothing about the shape of the law, but they give only a tail bound.
  \item \textbf{Dominating Gamma.} The dominating Gamma gives a law that propagates through the dynamics, but it commits to one family and adds some conservatism.
\end{itemize}

\section{SOS2 encodings and shortcuts that fail}

Reason: the code uses one selector binary per segment (encoding (i)), and the statement of encoding (i) stays in the scoping document.

The item ``Encodings'' of the scoping document lists two encodings of the union of tangents.
Encoding (i) stays there.
The removed part is encoding (ii), the statement that the two encodings describe the same set, and the shortcuts that fail.

\begin{enumerate}[nosep,label=(\roman*),start=2]
      \item An auxiliary variable $y$, tied by a piecewise-linear \emph{equality} (the SOS2 $\lambda$-formulation) to the pointwise minimum of the tangents to $\sqrt{v}$ at
        \begin{equation}
          v_r \;=\; (\tau - \mu_r)^2/\kappa_C^2,
        \end{equation}
        followed by the single linear row $\mu + \kappa_C\, y \le \tau$.
\end{enumerate}

    The two encodings describe the same set: the tangent at $\mu_r$ on the side of~$\mu$ and the tangent at $v_r$ on the side of~$v$ bound the same half-plane
    \begin{equation}
      \kappa_C^2 v \;\le\; (\tau - \mu_r)(\tau + \mu_r - 2\mu).
    \end{equation}

\begin{itemize}[nosep]
  \item \textbf{Shortcuts that fail.} Two natural shortcuts both fail.
    Requiring $y \ge$ every tangent gives their pointwise \emph{maximum}: a region strictly inside the single-tangent region, worse than no refinement.
    Requiring $y \le$ every tangent lets $y$ fall below $\sqrt{v}$, and this destroys the certificate.
    Likewise, the equality must interpolate the vertices of the tangent envelope, not points on the curve.
    Consecutive tangents meet at
    \begin{equation}
      v \;=\; \sqrt{v_r v_{r+1}}
      \qquad\text{with value}\qquad
      \bigl(\sqrt{v_r} + \sqrt{v_{r+1}}\bigr)/2,
    \end{equation}
    above the curve.
    The chords between points on the curve underestimate the concave square root.
\end{itemize}

\section{Inverse moments of the scatter families; remark on Miner tails}

Reason: they are used only by interpretation (I), which the project does not use.

\paragraph{Inverse moments of the scatter families.}
Every quantity that the framework uses later is a moment of $1/N_f$.
For the three families of Table~\ref{S:tab:sn_scatter} and $r > 0$,
\begin{equation}\label{eq:miner_invmom}
  \mathbb{E}\big[N_f^{-r}\big] \;=\;
  \begin{cases}
    \exp\big({-r\mu + r^2\sigma^2/2}\big), & \ln N_f \sim \mathcal{N}(\mu,\sigma^2) \quad \text{(finite for all $r$)},\\[3pt]
    \lambda^{-r}\,\Gamma(1 - r/k), & N_f \sim \mathrm{Wei}(k,\lambda) \quad \text{(finite if and only if $r < k$)},\\[3pt]
    e^{-r\mu}\,\Gamma(1 - r\beta), & \ln N_f \sim \mathrm{Gumbel}_{\min}(\mu,\beta) \quad \text{(finite if and only if $r\beta < 1$)}.
  \end{cases}
\end{equation}
The derivation is in Appendix~\ref{drv:miner_invmom}, page~\pageref{drv:miner_invmom}.
In particular, for the three families,
\begin{equation}
  \operatorname{Var}[1/N_f] \;=\;
  \begin{cases}
    e^{-2\mu+\sigma^2}(e^{\sigma^2}-1), & \text{log-normal},\\[3pt]
    \lambda^{-2}\big[\Gamma(1-2/k)-\Gamma(1-1/k)^2\big], & \text{Weibull},\\[3pt]
    e^{-2\mu}\big[\Gamma(1-2\beta)-\Gamma(1-\beta)^2\big], & \text{Gumbel-minimum}.
  \end{cases}
\end{equation}
These values match the increment moments already given for the log-normal below Table~\ref{S:tab:sn_scatter}.

\begin{remark}[Two tail behaviours, one real restriction]\label{rem:miner_tails}
The Gumbel-minimum log-life is \emph{exactly} a Weibull life:
\begin{equation}
  \Pr(N_f > x) \;=\; \exp\big(-e^{(\ln x - \mu)/\beta}\big) \;=\; \exp\big(-(x/e^{\mu})^{1/\beta}\big),
\end{equation}
that is, $\mathrm{Wei}(k = 1/\beta,\ \lambda = e^{\mu})$.
The existence conditions in~\eqref{eq:miner_invmom} correspond to each other.
So the three families of Table~\ref{S:tab:sn_scatter} give two tail behaviours for $1/N_f$: log-normal (all polynomial moments finite) and Fr\'echet tails with index $k$ (moments finite only below order $k$).

The second behaviour is a real restriction, not a formality.
The mean of the increment requires $k > 1$, and the variance, the input of the Cantelli certificate~\eqref{S:eq:cantelli_surrogate}, requires $k > 2$.
So the fitted shape must be checked before the method is used.
For $k \in (1,2]$ only the Markov bound remains, and for $k \le 1$ the untruncated model has an infinite mean increment and cannot enter the framework at all.

Moreover, for \emph{all three} untruncated families, $1/N_f$ has unbounded support and an infinite moment generating function.
So the exponential certificates (Chernoff~\eqref{S:eq:chernoff_permission}, Bernstein~\eqref{S:eq:bernstein_constraint}, Hoeffding) are not admissible.
Only a lower bound $N_{\min} > 0$ on the life, defended on physical grounds, restores them (a truncated family).
This bound gives the per-mission support $b_{\ell j} = n_j / N_{\min}$, and the moments must then be recomputed under the truncated law (closed forms with Gaussian tails and incomplete gamma functions).
\end{remark}

\section{Interpretation (I) of the random fatigue life}

Reason: the case study and the code use interpretation (II).

\emph{(I) Block-independent scatter.}
Interpretation~(I) suits the case in which the variability from cycle to cycle and from day to day dominates, and the component does not carry one fixed $N_f$.
Assume a single dominant stress level per mission.
(For a spectrum in bins, the means add over the bins; the variance depends on the assumed dependence within the mission, and it is computed from~\eqref{eq:miner_invmom} in the same way.)
Then the per-mission inputs are
\begin{equation}\label{eq:miner_routeI}
  \mu_{\ell j} \;=\; n_j\,\mathbb{E}\big[N_f^{-1}\big],
  \qquad
  v_{\ell j} \;=\; n_j^2\,\operatorname{Var}\big[N_f^{-1}\big].
\end{equation}
Independence across days is exactly the hypothesis of Equation~\eqref{S:eq:additive_functionals}.
So the Cantelli certificate~\eqref{S:eq:cantelli_surrogate} and the moment loop~\eqref{S:eq:loop_moments} apply \emph{without change}, with parametric instead of empirical moments per mission.
The rainflow route of Section~\ref{S:sec:cons_surrogates} already uses this interpretation.
The random-$N_f$ model is then only another input pipeline for the same formulation, subject to the admissibility conditions of Remark~\ref{rem:miner_tails}.

\section{Implications for the abstract formulation}

Reason: a deterministic-equivalent model type is not implemented and not planned.

\paragraph{Implications for the abstract formulation.}
\begin{itemize}[nosep]
  \item \textbf{Interpretation~(I).} Interpretation~(I) requires no change to Problem~\eqref{S:eq:abstract}.
    It supplies the per-mission moments~\eqref{eq:miner_routeI} to the same Route-I pipeline that uses the moments measured by rainflow counting, and it keeps the certificates and the loop condition of that pipeline.
  \item \textbf{Interpretation~(II).} Interpretation~(II) defines a strictly simpler model family (Table~\ref{S:tab:miner_interpretations}).
    There is no moment tracking, no slack from a tail bound and no big-M linearisation of a variance recursion.
    One quantile $q_\varepsilon$ per component holds the entire stochastic description.
    So interpretation~(II) is a natural candidate for a dedicated model type (a deterministic equivalent) beside the models based on processes.
\end{itemize}
The caveat of Remark~\ref{S:rem:miner_dichotomy} applies: which interpretation matches a given component is an empirical question, not a preference of the formulation.

\section{Compound Poisson rows of the three ordering tables}

Reason: compound Poisson is not implemented.

\noindent\textit{Row of Table~\ref{S:tab:ordering_lr} (likelihood ratio order).}

\begin{center}
\small
\begin{tabular}{@{}p{2.1cm}p{3.8cm}p{4.0cm}p{4.0cm}@{}}
\toprule
\textbf{Family} & $X$ & $Y$ & \textbf{Condition} \\
\midrule
Comp.\ Poisson
  & $\mathrm{CP}(\nu_1,\,F_{Y_1})$
  & $\mathrm{CP}(\nu_2,\,F_{Y_2})$
  & --- (can fail even for $F_{Y_1}{=}F_{Y_2}$, $\nu_1 \le \nu_2$; holds, for example, for deterministic jumps; see text) \\
\bottomrule
\end{tabular}
\end{center}

\noindent\textit{Row of Table~\ref{S:tab:ordering_hr} (hazard rate order).}

\begin{center}
\small
\begin{tabular}{@{}p{2.1cm}p{3.8cm}p{4.0cm}p{4.0cm}@{}}
\toprule
\textbf{Family} & $X$ & $Y$ & \textbf{Condition} \\
\midrule
Comp.\ Poisson
  & $\mathrm{CP}(\nu_1,\,F_{Y_1})$
  & $\mathrm{CP}(\nu_2,\,F_{Y_2})$
  & --- (no general criterion; the $\le_{\mathrm{lr}}$ route is unavailable, see text) \\
\bottomrule
\end{tabular}
\end{center}

\noindent\textit{Row of Table~\ref{S:tab:ordering_st} (usual stochastic order).}

\begin{center}
\small
\begin{tabular}{@{}p{2.1cm}p{3.8cm}p{4.0cm}p{4.0cm}@{}}
\toprule
\textbf{Family} & $X$ & $Y$ & \textbf{Condition} \\
\midrule
Comp.\ Poisson
  & $\mathrm{CP}(\nu_1,\,F_{Y_1})$
  & $\mathrm{CP}(\nu_2,\,F_{Y_2})$
  & for $F_{Y_1} = F_{Y_2}$, $Y \ge 0$, $\Pr(Y > 0) > 0$: $\nu_1 \le \nu_2$ \\
\bottomrule
\end{tabular}
\end{center}

\section{Paragraph ``Compound Poisson'' of the stochastic orderings}

Reason: compound Poisson is not implemented.

\paragraph{Compound Poisson.}
Take the same jump-size distribution~$F_Y$ for both laws, with $Y \ge 0$ and $\Pr(Y>0) > 0$.
Then
\begin{equation}
  \mathrm{CP}(\nu_1, F_Y) \le_{\mathrm{st}} \mathrm{CP}(\nu_2, F_Y)
\end{equation}
if and only if $\nu_1 \le \nu_2$.
The derivation is in Appendix~\ref{drv:order_cp}, page~\pageref{drv:order_cp}.

The stronger orders do \emph{not} follow in general.
Take $Y$ uniform on $\{1, 10\}$, $\nu_1 = 0.1$ and $\nu_2 = 0.2$.
The likelihood ratio drops from ${\approx}463$ at $x = 9$ (reachable only by nine unit jumps) to ${\approx}1.8$ at $x = 10$ (reachable by a single large jump), so $\le_{\mathrm{lr}}$ fails.
The ratio of the survival functions is not monotone either, so $\le_{\mathrm{hr}}$ fails as well.
Both orders hold in special cases, such as deterministic jump sizes, where the process is a scaled Poisson process.

When the jump-size distributions differ, the ordering conditions depend on the specific form of~$F_Y$, and no universal closed-form characterisation is available.

\section{Paragraph ``Compound Poisson'' of the loop constraint}

Reason: compound Poisson is not implemented.

\paragraph{Compound Poisson.}
When $D \sim \mathrm{CP}(\nu,F_Y)$, the survival function does not admit a universal closed form.
Conditioning on the number of shocks gives the series
\begin{equation}\label{eq:surv_cp}
  \Pr(D > \tau) \;=\; \sum_{n=0}^{\infty} \frac{e^{-\nu}\,\nu^n}{n!}\;\Pr\!\Bigl(\sum_{k=1}^{n} Y_k > \tau\Bigr),
\end{equation}
which can be evaluated numerically (it converges rapidly) or bounded.

\section{Upper bounds when no closed form is available}

Reason: the paragraph repeats the conservative surrogates of Section~\ref{S:sec:cons_surrogates}.

\paragraph{Upper bounds when no closed form is available.}
Sometimes the survival function $\Pr(D > \tau)$ has no closed form.
This happens for the compound Poisson process, and when the closure condition fails and the marginal of the sum is not in a standard family.
The \emph{reliability} constraint can then still be enforced through conservative upper bounds on the failure probability:
\begin{itemize}[nosep]
  \item Markov ($\Pr(D > \tau) \le \mathbb{E}[D]/\tau$);
  \item Cantelli~\eqref{S:eq:cantelli_surrogate};
  \item Chernoff~\eqref{S:eq:chernoff_surrogate}.
\end{itemize}
Section~\ref{S:sec:cons_surrogates} develops these bounds in per-mission form, together with the Bernstein and Hoeffding refinements, their relative tightness in the small-$\varepsilon$ regime and their MILP encodings.
For the compound Poisson process the Chernoff bound is particularly convenient, since $M_D(s) = \exp\bigl(\nu\,(M_Y(s)-1)\bigr)$ has a closed form whenever $M_Y$ has one.


\section{Maintenance under the Palmgren--Miner damage model}

Reason: there is no separate Miner maintenance model; the rainflow repair maps cover the code.

\paragraph{Maintenance under the Palmgren--Miner damage model.}
The remaining-life framework of Section~\ref{S:sec:remaining_life} represents damage as the normalised fraction
\begin{equation}
  D_{\mathrm{PM}} = \sum_i n_i/N_{f,i} \ge 0,
\end{equation}
with failure at $D_{\mathrm{PM}} \ge 1$.
$D_{\mathrm{PM}}$ adds up over loading blocks.
So it fits the same damage-tracking structure as the stochastic processes of Table~\ref{S:tab:deg_models}: each mission day adds an increment $\Delta D_{\mathrm{PM}} = n_j / N_{f,j}$.
This increment is deterministic if the S--N parameters are fixed, and random if the scatter is modelled~\cite{MunizCalvente2022}.
The repair operator~$\mathcal{R}$ and the replacement therefore apply directly:
\begin{itemize}
  \item \emph{Replacement} resets the fatigue fraction to zero: $D_{\mathrm{PM}}^+ = 0$ (equivalently $D^{\mathrm{new}} = 0$).
    A new component, with a cycle count of zero, takes the place of the physical component.
  \item $\text{ARD}_\infty$:
    \begin{equation}
      D_{\mathrm{PM}}^+ = (1-\rho)\,D_{\mathrm{PM}}^-.
    \end{equation}
    Physically, this models an intervention that partly restores the fatigue life without resetting the whole loading history, for example a surface treatment, shot peening or a local crack repair.
  \item $\text{ARD}_1$:
    \begin{equation}
      D_{\mathrm{PM}}^+ = D_{\mathrm{PM}}^- - \rho\,(D_{\mathrm{PM}}^- - D_{\mathrm{PM}}^{\mathrm{last}}).
    \end{equation}
    The repair reverses part of the damage accumulated since the last intervention (repair or replacement), and only that damage.
\end{itemize}
Unlike the stochastic process models, $D_{\mathrm{PM}}$ is normalised so that failure occurs at~$1$.
So the threshold is fixed at $\tau = 1$, and the reliability constraint becomes
\begin{equation}
  \Pr(D_{\mathrm{PM}} \ge 1) \le \varepsilon.
\end{equation}
In the deterministic variant of Miner's rule this is a hard feasibility check.
In the stochastic extension it becomes a chance constraint analogous to~\eqref{S:eq:abs_rel}.

Closure under summation holds trivially: $D_{\mathrm{PM}}$ is a sum of real-valued increments.
When the increments are independent, as under the block-independent scatter~(I) of Section~\ref{S:sec:miner_randomN}, the marginal on any day is the convolution of the individual increments, whatever their distributional family.
Under the unit-specific scatter~(II) the increments share the factor~$A$, and the marginal is the law of $M_k/A$.
\begin{itemize}[nosep]
  \item In the deterministic variant of Miner's rule, the comparison of two schedules reduces to the comparison of two real numbers, and no stochastic order is needed.
  \item In the stochastic extension (scatter in $N_{f,i}$), the parametric ordering conditions of Section~\ref{S:sec:stoch_orders} apply only if the assumption on the scatter places the increments $n_i/N_{f,i}$ in one of the families of Table~\ref{S:tab:deg_models}.
    Otherwise the general definitions of the orders still hold, but they must be verified case by case.
\end{itemize}

\section{Brake friction elements, wheel bearings and traction-motor insulation}

Reason: these components are not tracked in the case study and not planned.

\paragraph{Brake friction elements.}
Friction braking wears the brake pad and the brake disc.
\begin{itemize}[nosep]
  \item \textbf{Model.}
    A physics-based route applies Archard's law on a fractal, elastoplastic contact~\cite[Eqs.~11,~14]{Sha2022}.
    Sha et al.\ apply it to the brake disc, not to the pad.
    The model gives the average wear depth $h_{\mathrm{wear}}$ of the disc in one braking~\cite[Eq.~21]{Sha2022}.
    The life of the disc, as a number of brakings, is then
    \begin{equation}
      N = H_0/(S_0\,h_{\mathrm{wear}}),
    \end{equation}
    with $H_0$ the maximum admissible average wear depth and $S_0$ a dimensionless wear safety factor for frictional heat and other external effects~\cite[Eq.~22]{Sha2022}.
  \item \textbf{Regenerative braking.}
    In an electric vehicle the regenerative brake absorbs most service braking, and it cuts friction-brake wear by an estimated 60--90\%~\cite[Sec.~6.1]{Fussell2022}.
    The remaining increments concentrate in the events where regeneration is unavailable or insufficient: a battery near full charge, a very low temperature, an emergency stop.
  \item \textbf{Fit to the framework.}
    The cumulative wear of a friction element remains a sum of non-negative, mission-dependent increments, one per friction-braking event.
    So it is a monotone accumulation process, whose rate now also depends on the energy-management strategy.
\end{itemize}

\paragraph{Wheel bearings.}
Bearing degradation has two stages: a long \emph{initiation} phase of rolling-contact fatigue, and then a rapid \emph{spall propagation} once a surface defect forms (a spall is a flake of material that breaks off the surface).
\begin{itemize}[nosep]
  \item \textbf{Model.}
    Ohana et al.~\cite[Sec.~2.1.3]{Ohana2023} model the propagation phase with a formulation from continuum damage mechanics.
    In it, a scalar damage variable accumulates per over-rolling cycle,
    \begin{equation}
      \mathrm{d}D/\mathrm{d}N = \bigl(\sigma_{\mathrm{eff}}/[\sigma_r(1-D)]\bigr)^m.
    \end{equation}
    The growth is driven by the tensile residual stresses that repeated impacts of the rolling elements generate at the edge of the spall.
    The model is meant as the basis of a physics-based prognostic tool that estimates the remaining useful life once a critical spall size is reached.
  \item \textbf{Fit to the framework.}
    The underlying driver is the number of load cycles, and hence the running distance and the speed.
    So the bearing damage is a monotone accumulation process.
    A Gamma or inverse Gaussian process (initiation) that feeds a remaining-life threshold (propagation) represents it well.
\end{itemize}

\paragraph{Traction-motor insulation.}
The winding insulation is the main element that limits the life of the traction motor, and its dominant ageing mechanism is thermal.
A high winding temperature accelerates the chemical degradation of the dielectric.
\begin{itemize}[nosep]
  \item \textbf{Model.}
    Dakin recognised thermal ageing as a chemical reaction whose rate constant follows the Arrhenius exponential law $K = A\,e^{-B/T}$ in absolute temperature~\cite[Eq.~5]{Dakin1948}.
    This gives the thermal-endurance life $L(T) = a\,\exp(b/T)$~\cite[Eq.~1]{Zhang2025insul}.
  \item \textbf{Montsinger's rule.}
    The earlier empirical rule of Montsinger states that the ageing rate roughly doubles for each 8\,$^\circ$C rise of temperature~\cite[Synopsis, p.~776]{Montsinger1930}.
    This rule halves the life for each fixed rise, and Dakin showed that it is the constant-coefficient approximation of the Arrhenius law, accurate at one temperature only~\cite[Fig.~1 and discussion of Eq.~5]{Dakin1948}.
  \item \textbf{Fit to the framework.}
    Under a temperature profile that varies with the mission, the fraction of life consumed in an interval $\mathrm{d}t$ at temperature $T_i(t)$ is $\mathrm{d}t/L[T_i(t)]$.
    The total consumed fraction is the integral of this loss rate over the profile.
    This is a cumulative-damage (Palmgren--Miner) extension of the Arrhenius model~\cite[Eq.~2 and Fig.~10]{Zhang2025insul}.
    So the abstract description coincides with the remaining-life model of Section~\ref{S:sec:remaining_life}, with the mission driving the thermal load (current, cooling, ambient conditions).
  \item \textbf{Stochastic processes.}
    The same review also surveys Wiener- and Gamma-process formulations of insulation degradation~\cite[Sec.~5]{Zhang2025insul}, which match directly the stochastic families of Section~\ref{S:sec:stoch_orders}.
\end{itemize}

\section{Input example of the specification}

Reason: the specification keeps the example, and this text only repeats it.

The input example of the specification has $F = 3$ members, $M = 2$ missions and $L = 2$ components.
Its six components use the Gaussian, inverse Gaussian, Wiener and Gamma models, and the rainflow model with the Bernstein bound and with the Chernoff bound.
It passes every input check.
\begin{itemize}[nosep]
  \item \textbf{Inverse Gaussian component.}
    It uses $\text{ARD}_1$ with the starting anchor
    \begin{equation}
      0.06 \le \mu_0 = 0.10.
    \end{equation}
  \item \textbf{Gamma component.}
    It has the starting anchor
    \begin{equation}
      0.05 \le \mu_0 = 0.08
    \end{equation}
    and no rate for the initial damage, so its initial damage is known.
  \item \textbf{Chernoff component.}
    It has
    \begin{equation}
      K_0 = 5 \cdot 0.14 = 0.70
    \end{equation}
    by default, so its starting anchor $0.50$ is allowed.
  \item \textbf{Bernstein component.}
    It has the support constant
    \begin{equation}
      b = \max(0.05, 0.03, 0.05) = 0.05,
    \end{equation}
    so
    \begin{equation}
      2\kappa b/3 = 0.077 < \tau = 1,
    \end{equation}
    with $\kappa = \ln 10$.
\end{itemize}

\section{Changelog of the software specification}
\label{sec:spec_changelog}

Reason: history of the specification (decision of 10 November 2026).


This section records the versions of the software specification \texttt{spec/spec.tex}.
Versions 0.3 to 0.5 use older names and an older horizon.
Their law $\text{ARD}_1$ is the law $\text{ARD}_\infty$ of this document, and their law $\text{ARA}_1$ is the law $\text{ARD}_1$ of this document.
Their horizon has $2H$ days, and the second $H$ days are the schedule to repeat.

\paragraph{Version 0.3 (panel review).}
\begin{itemize}[nosep]
  \item The $\text{ARA}_1$ maintenance was removed from the Gaussian model, because no source in the literature supports its extension to the variance.
  \item The output \texttt{v} has the shape $F \times L \times 2H$, with NaN (not a number) at the components without a variance.
  \item The bound on $\varepsilon$ became tighter: from $(0, 0.5)$ to $(0, 0.01]$.
  \item A verified proposition, with a proof by the monotone likelihood ratio (MLR), replaced the open item on first-order stochastic dominance (FSD) of the inverse Gaussian (IG) model.
  \item The remark on tightness became an argument with three mechanisms.
\end{itemize}

\paragraph{Version 0.4 (rainflow extension).}
\begin{itemize}[nosep]
  \item The rainflow degradation model was added.
    It is distribution-free, and its inputs are the empirical mean, variance and support of each mission, from rainflow counting.
    A Cantelli or a Bernstein tail bound certifies the reliability; the new option \texttt{tail\_bound} selects the bound, with the default \texttt{"cantelli"}.
    The specification gives both an exact nonconvex quadratic formulation and a conservative linear tangent formulation.
    The loop constraint acts on both tracked moments.
    The argument is that the componentwise domination of the moments chains across repetitions, while the scalar comparison of the bounds, $U_{2H} \le U_H$, does not.
    The model accepts $\text{ARD}_1$ only, as the Gaussian model did.
  \item The value \texttt{"socp"} of the key \texttt{formulation} was renamed \texttt{"exact"}.
    \texttt{"socp"} remained as a deprecated alias, with a \texttt{UserWarning}.
  \item The exact reliability constraint of the Gaussian and Wiener models had been called a rotated second-order cone (SOC); this was corrected.
    The variance is affine in the decisions, so the squared form is a nonconvex quadratic constraint, handled by the \texttt{NonConvex=2} support of Gurobi.
    The \texttt{"quadratic"} damage penalty remains a truly convex SOC constraint.
  \item The wording of the Wiener loop constraint was corrected after the accepted review item A-02: the constraint bounds the moments; it is not first-order stochastic dominance.
  \item The change log of version 0.3 above describes the revision that the title page had called ``v0.2 (revised)''.
  \item A panel review after the draft (markers RF-042--RF-045 and corrections made in place) made further changes.
    \begin{itemize}[nosep]
      \item It fixed the base case of the induction over the loop ($m \ge 2$, with the loop row as the case $m = 1$).
      \item It fixed the wording at the boundary of the precondition $\tau > 2\kappa b/3$.
      \item It fixed the centring step in the application of Bernstein's inequality, which holds a fortiori (it follows from a stronger fact).
      \item It fixed the statement of the problem class for fleets whose components are all inverse Gaussian or Gamma.
      \item It extended the consistency checks based on the Bhatia--Davis inequality (a bound on the variance from the mean and the range) to the initial and replacement states.
      \item It extended the list of checked array shapes to $b^{\mathrm{inc}}$, $b^{0}$ and $b^{\mathrm{new}}$.
      \item It moved \texttt{tail\_bound} to the input tables of the models.
      \item It added the fifth border style of the plot.
      \item It corrected the value of $\varepsilon$ in the YAML example, which did not respect the bound $(0, 0.01]$ of the specification.
    \end{itemize}
\end{itemize}

\paragraph{Version 0.5 (removal of the loop penalty).}
\begin{itemize}[nosep]
  \item The soft penalty of the loop constraint, $\mathcal{L}_{\mathrm{loop}}$, and its coefficient $C_P$ were removed from every place: the objective, the penalty formulas of the models, the total penalty sum, the input tables, the YAML example and the statement of the full problem.
    The hard loop constraints of all five models did not change.
  \item The remark on the tightness of the big-M rows had named three mechanisms.
    It was revised to name the two cost mechanisms that remain, $C_D\,u$ and $C_R\,z$.
  \item An argument on robustness to slack was added.
    The tracked states are componentwise upper bounds on the true moments, and each one-sided dynamics row dominates the exact monotone map of one day.
    So the hard loop constraint gives $\Pi(s_H) \le s_H$ for the exact map $\Pi$ of one repetition of the second half (Section~\ref{S:sec:perpetual}; that version wrote $\Phi$ for this map).
    This inequality is the premise of the induction in the repeatability argument.
    The monotonicity of the reliability certificates then keeps every future repetition certified, even when the tracked states carry slack.
    The only role of the penalty had been to break ties in favour of larger sustainability margins.
  \item A later review (markers REV-001--REV-022) closed a gap of the loop constraint under $\text{ARA}_1$: the loop did not cover the anchors.
    It added the hard rows
    \begin{align}
      \mu^{\mathrm{last}}_{2H} &\le \mu^{\mathrm{last}}_{H} && \text{(Wiener with $\text{ARA}_1$)}, \\
      \alpha^{\mathrm{last}}_{2H} &\le \alpha^{\mathrm{last}}_{H} && \text{(Gamma with $\text{ARA}_1$)}.
    \end{align}
    A derivation of monotonicity on the full state shows that these rows restore the chaining premise $\Pi(s_H) \le s_H$.
    The review extended the remark on robustness to slack in the same way.
  \item The same review made further changes.
    \begin{itemize}[nosep]
      \item It stated that an idle day is a row of zeros in the assignment, so the term with $C_M$ is a real decision cost, not a constant.
      \item It corrected the day index of the anchor in the $\text{ARA}_1$ repair-cost row ($k \to k{-}1$).
      \item It removed the old duplicate parameter \texttt{alpha\_target}, a former name of $\tau$.
      \item It made consistent which keys are optional: \texttt{tail\_bound} and the input fields that have a default.
      \item It added the derived constants to the notation table.
      \item It expanded the two items on future work of the rainflow model, with pointers to the Route-I refinements of this document (Section~\ref{S:sec:routeI_refinements}).
      \item It recorded two questions on the convention of the $\text{ARA}_1$ anchor as open items: whether the anchor stores the state before or after the repair, and the missing reset of the anchor at a replacement.
    \end{itemize}
\end{itemize}

\paragraph{Version 1.0 (merged package).}
\begin{itemize}[nosep]
  \item The specification describes the package that merges the branches \texttt{lore} and \texttt{main}, on the code base of \texttt{main}.
    It holds only the codebase and its mathematical constructions; the reasons, the proofs and the sources are in this scoping document.
  \item The five degradation models are cell builders in one registry; a cell is component~$\ell$ of member~$i$.
  \item The repair laws are $\text{ARD}_\infty$, which reduces the whole degradation, and $\text{ARD}_1$, which reduces the degradation since the last intervention.
  \item The horizon has a transitory phase and an operating phase, $H = [H_1, H_2]$, and the loop row bounds the full tracked state, anchors included.
  \item Every dynamics row is one-sided, in big-M or indicator form.
  \item The rainflow model offers five tail bounds, and the quadratic reliability rows have an exact, a tangent and a piecewise-linear form.
  \item The inverse Gaussian row is a cap from a conditional tail mean.
    The Gamma row is a cap on the tracked mean that holds for any mix of rates.
  \item The reliability level, the tangent point, the repair cost and the replacement cost are set per cell.
  \item The specification covers the horizon sweeps, the warm starts, the schedule plot and the validators.
  \item The input keys follow one table, and unknown or old keys raise an error.
  \item The input checks give every non-strict inequality the relative tolerance $10^{-12}\max(1, |a|, |b|)$, and they compare selector values with their case.
    At a key without a default, \texttt{null} means that the key is not given at that cell.
    YAML and JSON result files write every number that is not finite as \texttt{null}.
  \item The specification fixes the numerical methods of the two caps: a root search in $\ln q$ for the inverse Gaussian quantile, and the grid, the golden-section search and the two bisections of the Gamma cap~$A_1$.
  \item One function, \texttt{read\_result}, reads the result files for every tool.
    The greedy heuristic tries a repair and a replacement on every free day of a Wiener cell.
\end{itemize}

\section{Derivation of the inverse moments}

Reason: its statement, the inverse moments, is removed with interpretation (I).

\phantomsection\label{drv:miner_invmom}
\begin{proof}[Derivation of the inverse moments~\eqref{eq:miner_invmom}]\leavevmode
\begin{itemize}[nosep]
  \item \textbf{Log-normal.} The first line is the normal moment generating function at the argument $-r$.
  \item \textbf{Weibull.} The second line follows from the substitution $u = (x/\lambda)^k$, which turns $\mathbb{E}[N_f^{-r}]$ into
    \begin{equation}
      \lambda^{-r}\int_0^\infty u^{-r/k}e^{-u}\,\mathrm{d}u;
    \end{equation}
    this integral diverges \emph{at the origin} when $r \ge k$.
  \item \textbf{Gumbel-minimum.} The third line writes $\ln N_f = \mu + \beta G$ with $\Pr(G \le g) = 1 - e^{-e^{g}}$.
    The transform
    \begin{equation}
      \mathbb{E}[e^{tG}] \;=\; \int_0^\infty u^{t}e^{-u}\,\mathrm{d}u \;=\; \Gamma(1+t)
    \end{equation}
    (substitute $u = e^g$) is finite if and only if $t > -1$.
\end{itemize}
\end{proof}

\section{Derivation of the compound Poisson row of the ordering table}

Reason: its statement, the compound Poisson row, is removed.

\phantomsection\label{drv:order_cp}
\begin{proof}[Derivation of the compound Poisson row of Table~\ref{S:tab:ordering_st}]
Sufficiency follows from the superposition coupling
\begin{equation}
  \mathrm{CP}(\nu_2, F_Y) \stackrel{d}{=} \mathrm{CP}(\nu_1, F_Y) + \mathrm{CP}(\nu_2 - \nu_1, F_Y),
\end{equation}
with independent non-negative summands.
Necessity follows from the atom at zero, $\Pr(X = 0) = e^{-\nu \Pr(Y > 0)}$.
\end{proof}

\bibliographystyle{plainnat}
\bibliography{scoping_bib}

\end{document}
